\subsection{Sprint 2 --- 15/09/2023 até 06/10/2023}

Após a retrospectiva da Sprint 1, a equipe iniciou uma nova fase buscando aplicar as melhorias identificadas e manter o alinhamento com os objetivos do projeto. A Sprint 2 representou uma continuação estratégica do desenvolvimento, mas também trouxe desafios mais intensos, especialmente relacionados ao backend e à gestão de tempo.

Nesta sprint, foram definidas duas \textit{user stories} principais. A primeira envolvia a visualização de informações da sessão, enquanto a segunda tratava da importação de arquivos \texttt{.csv}. Diferentemente da sprint anterior, a ênfase desta etapa recaiu mais sobre o backend do que sobre o frontend, o que representou um desafio significativo para mim.

Durante a sprint, enfrentei dificuldades para contribuir de forma mais consistente. O primeiro motivo foi minha pouca experiência com backend, especialmente com Node.js. O segundo foi a sobrecarga de provas e trabalhos no meio do semestre, que tornou mais difícil administrar o tempo dedicado às atividades da AGES.

Mesmo com essas limitações, consegui desenvolver uma familiaridade inicial com Node.js, aprendizado que se mostrou importante para as sprints seguintes. Minhas contribuições foram mais limitadas do que eu gostaria, mas a sprint ainda teve valor formativo por me expor a uma frente técnica diferente e por evidenciar a necessidade de ampliar minha atuação além do frontend.

As dificuldades de gestão de tempo não foram exclusivas minhas. A equipe como um todo enfrentou problemas nesse aspecto, o que impactou o projeto. Após o bom desempenho da Sprint 1, percebi que houve certo relaxamento coletivo, resultando em contratempos. Embora esses problemas não tenham comprometido de forma definitiva o projeto, funcionaram como um alerta importante sobre a necessidade de disciplina e acompanhamento contínuo.

O encerramento da sprint ocorreu na terceira reunião com os stakeholders. Apesar das dívidas técnicas existentes, a reunião foi tranquila. Os membros mais experientes conseguiram transmitir confiança aos stakeholders, demonstrando que a equipe tinha condições de contornar os desafios enfrentados. Também ficou claro que esse período do semestre costuma ser naturalmente mais difícil em projetos AGES, devido à concentração de demandas acadêmicas.

Em síntese, a Sprint 2 foi marcada por desafios de tempo, adaptação técnica e contribuição limitada da minha parte. Ainda assim, considero que ela teve papel importante no meu aprendizado, pois mostrou a necessidade de organização mais rigorosa e abriu caminho para minha aproximação posterior com o backend.
